\documentclass[../elsevier.tex]{subfiles}
\begin{document}

\section{The Camouflage-CSBM}
\label{sec:model}

We work in a two-class contextual stochastic block model augmented with an
adversarial parameter controlling how aggressively bots attach to humans.

\begin{definition}[Camouflage-CSBM]
\label{def:ccsbm}
Let $\mathcal{V}$ be a set of $N$ nodes partitioned into bots $\mathcal{B}$ and
humans $\mathcal{H}$, with bot prevalence $\pi_B = |\mathcal{B}|/N \in (0,1)$ and
bot-to-human ratio $c = \pi_B/(1-\pi_B)$. Each node $v$ carries a feature vector
\begin{equation}
    \mathbf{x}_v \sim \mathcal{N}(\boldsymbol{\mu}_{y_v},\, \sigma^2 I_d),
    \qquad y_v \in \{\mathrm{B}, \mathrm{H}\},
\end{equation}
and we write $\Delta = \boldsymbol{\mu}_{\mathrm{B}} - \boldsymbol{\mu}_{\mathrm{H}}$
for the class-mean gap. Every node has neighbourhood size $D$. Neighbours are drawn
independently: a bot's neighbour is human with probability $\rho$ (the
\emph{camouflage ratio}) and a bot otherwise; a human's neighbour is a bot with
probability $\eta$ and a human otherwise.
\end{definition}

The two mixing parameters are not free. Counting bot--human edge endpoints from
both sides gives $|\mathcal{B}| D \rho = |\mathcal{H}| D \eta$, hence the
\emph{edge-balance identity}
\begin{equation}
    \eta(\rho) \;=\; \frac{\pi_B}{1-\pi_B}\,\rho \;=\; c\,\rho .
    \label{eq:balance}
\end{equation}
Camouflage is therefore doubly harmful: it not only pollutes bot neighbourhoods but
necessarily pollutes human neighbourhoods as well, at a rate set by the bot
prevalence. Every result below accounts for this coupling; treating $\eta$ as
independent of $\rho$ overstates a detector's robustness.

\begin{assumption}[Regularity used by the depth analysis]
\label{as:tree}
Definition~\ref{def:ccsbm} samples each actor's alters independently, which suffices
for the single-layer results. Section~\ref{sec:depth} additionally requires that the
tie relation be \emph{symmetric}---so that $v \in \mathcal{N}(u)$ whenever
$u \in \mathcal{N}(v)$, which is what makes~\eqref{eq:backtrack} hold---and that the
neighbourhood be locally tree-like with branching $D$, so that the number of actors at
distance $j$ is $N_j = D(D-1)^{j-1}$. Both are standard for sparse networks in which
short cycles are rare, and neither is needed before Section~\ref{sec:depth}. We also
take $\rho \le 1/c$, so that $\eta \le 1$, which holds automatically whenever the
strategic class is the minority.
\end{assumption}

\begin{remark}[Modelling choices]
Independent neighbour sampling at fixed size $D$ makes the analysis exact and is
the discrete analogue of the degree-concentration assumptions used
in~\citep{baranwal2021graph}. Isotropic within-class covariance and a common $D$ are
simplifications; Section~\ref{sec:limitations} discusses what changes without them.
We emphasise that $\rho$ is an \emph{adversarial} parameter chosen by the attacker,
not a fixed property of the generative process---this is the substantive
difference from the standard CSBM setting.
\end{remark}

\paragraph{Aggregation operator.}
We analyse one round of mean aggregation, the canonical homophilic operator,
\begin{equation}
    \mathbf{h}_v \;=\; \frac{1}{D}\sum_{u \in \mathcal{N}(v)} \mathbf{x}_u ,
    \label{eq:meanagg}
\end{equation}
which isolates the effect of neighbourhood averaging from the confounds of learned
weights and nonlinearity. Because all the quantities we need are governed by the
projection onto the discriminant direction, define
$\mathbf{w} = \Delta/\|\Delta\|$ and $t_v = \mathbf{w}^\top \mathbf{h}_v$.

\paragraph{Baseline for comparison.}
The natural point of reference is the graph-free classifier, which uses
$\mathbf{x}_v$ directly. For two isotropic Gaussians with equal priors its Bayes
error is
\begin{equation}
    \varepsilon_{\mathrm{raw}} \;=\; \Phi\!\left(-\frac{\|\Delta\|}{2\sigma}\right),
    \label{eq:rawerr}
\end{equation}
with $\Phi$ the standard normal CDF. The question driving Section~\ref{sec:theory}
is when aggregation beats~\eqref{eq:rawerr} and when it does not.

\end{document}
